
We describe the architecture and ablations of \textbf{Qwen3.8-Flash-Next}, a sparse mixture-of-experts model with 125B parameters, 6B activated per token, and additional 51B parameters of n-gram embedding tables held off the accelerator.
On fourteen pre-training benchmarks the model leads the 397B-A17B predecessor on eight and trails it on the rest by at most 2.6 points, at \textbf{1/3} the activated parameters, \textbf{1/3} the training tokens, and roughly \textbf{1/9} the training FLOPs.
Token mixing uses a layer-wise hybrid of \textbf{Gated DeltaNet (GDN)} and global attention, with one full-attention layer in every four; at continued-pretraining time those full-attention layers are replaced by \textbf{Qwen Sparse Attention (QSA)}, which scores context at micro-block granularity with a compressed lightweight indexer. The residual stream is widened to four branches and read through an elementwise gate, a design we call the \textbf{Gated Residual (GR)}. 
Capacity is added outside the backbone by a single \textbf{n-gram embedding layer} whose tables are prefetched from host memory.
We evaluate every candidate change along three axes: loss together with downstream benchmarks; the cost of the change in training, prefill and decode; and its effect on the optimal hyperparameters and training stability. 
Loss and downstream accuracy do not always move together: enlarging the n-gram vocabulary lowers loss monotonically while downstream accuracy saturates.
The architecture and the \textbf{Muon} optimizer together shift the optimal learning rate and batch size upwards, render batch-size warmup unnecessary, and substantially improve stability under stress tests.
Loss, benchmarks, efficiency and stability form one design problem. 
Solved jointly, they yield a recipe that is simultaneously \textit{more efficient, more capable and more stable}.